\chapter{The Four Faces}
\label{chap:the-four-faces}

The three trips left a number in a bracket. This chapter reads it --- and the number turns out to have four
faces, which the machine reads at three different costs. That the costs differ is not an accident to be
smoothed over; it is the chapter's whole point. What a face of the number costs to know is what that face
\emph{is}. Two of the four are free, read straight off the bracket. One asks the machine to make a decision.
And the fourth is not read at all in the ordinary sense: it is measured, the machine forcing its own single
needle and reporting what the forcing cost. Four faces, three costs, and the stratification of the costs is
the epistemology of the whole construction, made visible on one number.

\section{The two free faces}
\label{se:two-free-faces}

Two faces of the number are already in hand, because the bracket put them there. The charge is the lower
wall of the bracket --- the loop count the seam accumulated --- and the mass is the upper wall, the
curvature the third trip recovered. To read either is to project: to take a bracket already built and name
one of its two walls, which costs nothing beyond the naming. These are the free faces. They are free because
the three trips already paid for them; the number arrived bracketed, and its two bounds are its charge and
its mass, there for the reading. The construction \emph{tanges} the charge from the mass --- they are
different walls, told apart by which end of the bracket they sit at --- but neither telling costs anything,
because both were fixed before this chapter began.

\section{The decision}
\label{se:the-decision}

The third face costs a decision. The phase of the number --- which of two orientations it carries --- is not
sitting in the bracket to be projected; it must be read by asking the machine to decide a fact, the carried
fact whose two ways the phase distinguishes. This is the local gauge, the orientation read at a point, and
reading it is not free: it costs exactly one act of decision, the machine forcing the carried fact to fall
one way or the other. The cost is small and it is definite --- a single decision, concentrated at the one
place the fact is decided --- but it is a cost, and it is a different kind of thing from projecting a wall of
the bracket. The phase is what the machine must \emph{ask} to know; the charge and mass are what it already
has. That difference in cost is a difference in kind, and the construction keeps it visible.

\section{The measured face}
\label{se:the-measured-face}

The fourth face is not read like the others; it is measured. It is the sameness --- the agreement that the
truth the machine carried down came back as the truth it set out with --- and this the construction does not
prove as a theorem but reads as a measurement. The machine forces its one needle, the single decision the
whole descent bottoms out on, and reports the cost of forcing it: the heartbeats the elaboration spends to
settle that one fact. That cost is a real, reproducible quantity --- the machine can read its own pulse, and
the reading is concentrated exactly at the needle, the one place the construction commits. This is the
deepest sense in which the number has a face the machine measures rather than derives: it puts its own act
of deciding on the scale and reads the weight.

One honesty must ride with this reading. The pulse is a real cost, but its absolute size is not an
invariant of the machine alone; it depends on the state the elaboration is in when the reading is taken ---
on what work has already been done and cached before the needle is forced. Measure the same needle in a
cold machine and a warm one and the number differs, not because the needle changed but because the reading
did. So the construction reports the pulse as a reading, never as a fixed constant: the robust fact is that
the machine \emph{can} weigh its own single decision, and that the weight sits at the needle and nowhere
else. The magnitude is a measurement with a context; the location of the cost is the invariant. This is the
most thoroughly measured, and the most carefully fenced, of the four faces --- a number the machine reads
off itself and reports honestly, caveat and all.

\section{The butterfly}
\label{se:the-butterfly}

One reading of the charge face is worth naming, offered as interpretation and not as a further theorem. The
charge counts the loop the way a certain bit-count counts a set: the positron count is the number of set
bits in one combination of the electron pattern and the number, and the matter reading is its complement, so
that matter and positron together make the whole population of set bits in the electron pattern. Read this
way the charge is a population count: the construction \emph{funges} the resolved places by which way they
fell, pooling them into two tallies --- how many came back one way, how many the other --- and the matter
and the antimatter are those two, partitioning the whole between them. It is the same currency
the second book counted annihilation in; here it is the charge face of the number, read as a count of which
way each resolved place fell. This is a reading laid over the proved charge, named as a reading.

\section{What is demonstrated, and what is assumed}
\label{se:demonstrated-assumed-III6}

What is demonstrated is the four faces and their three costs: the charge and the mass projected free off the
bracket, the phase read by one decision, and the sameness measured as the machine's own pulse at the needle.
The cost stratification is the demonstrated content --- two free, one decided, one measured --- and it is
the epistemology of the construction shown on a single number: what each face of the number costs to know is
exactly what that face is. The measured face is reported with its caveat: a real pulse, its magnitude
context-dependent, its cost located invariantly at the needle.

What is assumed is unchanged. The seed and its decidability; the honest import, marked as borrowed; the
single standing grant, the law of motion the first book named at its summit, still out of play. And one
naming is held back on purpose: which particular value of the one invariant this number is --- the identity
that would let it be called an electron and no other --- is reached for and named only at the close, not
here. This chapter reads the number's four faces; it does not yet name the number. That the pulse can itself
be predicted from its parts, and what the residual of that prediction means, is a further matter, taken up
when the machine's costs are weighed whole. With the number read four ways and its costs stratified, one
question remains for the movement: what the whole construction has produced, over its two carriers, taken
together. That is the next chapter.
